% Options for packages loaded elsewhere
% Options for packages loaded elsewhere
\PassOptionsToPackage{unicode}{hyperref}
\PassOptionsToPackage{hyphens}{url}
\PassOptionsToPackage{dvipsnames,svgnames,x11names}{xcolor}
%
\documentclass[
  letterpaper,
  DIV=11,
  numbers=noendperiod]{scrartcl}
\usepackage{xcolor}
\usepackage{amsmath,amssymb}
\setcounter{secnumdepth}{5}
\usepackage{iftex}
\ifPDFTeX
  \usepackage[T1]{fontenc}
  \usepackage[utf8]{inputenc}
  \usepackage{textcomp} % provide euro and other symbols
\else % if luatex or xetex
  \usepackage{unicode-math} % this also loads fontspec
  \defaultfontfeatures{Scale=MatchLowercase}
  \defaultfontfeatures[\rmfamily]{Ligatures=TeX,Scale=1}
\fi
\usepackage{lmodern}
\ifPDFTeX\else
  % xetex/luatex font selection
\fi
% Use upquote if available, for straight quotes in verbatim environments
\IfFileExists{upquote.sty}{\usepackage{upquote}}{}
\IfFileExists{microtype.sty}{% use microtype if available
  \usepackage[]{microtype}
  \UseMicrotypeSet[protrusion]{basicmath} % disable protrusion for tt fonts
}{}
\makeatletter
\@ifundefined{KOMAClassName}{% if non-KOMA class
  \IfFileExists{parskip.sty}{%
    \usepackage{parskip}
  }{% else
    \setlength{\parindent}{0pt}
    \setlength{\parskip}{6pt plus 2pt minus 1pt}}
}{% if KOMA class
  \KOMAoptions{parskip=half}}
\makeatother
% Make \paragraph and \subparagraph free-standing
\makeatletter
\ifx\paragraph\undefined\else
  \let\oldparagraph\paragraph
  \renewcommand{\paragraph}{
    \@ifstar
      \xxxParagraphStar
      \xxxParagraphNoStar
  }
  \newcommand{\xxxParagraphStar}[1]{\oldparagraph*{#1}\mbox{}}
  \newcommand{\xxxParagraphNoStar}[1]{\oldparagraph{#1}\mbox{}}
\fi
\ifx\subparagraph\undefined\else
  \let\oldsubparagraph\subparagraph
  \renewcommand{\subparagraph}{
    \@ifstar
      \xxxSubParagraphStar
      \xxxSubParagraphNoStar
  }
  \newcommand{\xxxSubParagraphStar}[1]{\oldsubparagraph*{#1}\mbox{}}
  \newcommand{\xxxSubParagraphNoStar}[1]{\oldsubparagraph{#1}\mbox{}}
\fi
\makeatother


\usepackage{longtable,booktabs,array}
\newcounter{none} % for unnumbered tables
\usepackage{calc} % for calculating minipage widths
% Correct order of tables after \paragraph or \subparagraph
\usepackage{etoolbox}
\makeatletter
\patchcmd\longtable{\par}{\if@noskipsec\mbox{}\fi\par}{}{}
\makeatother
% Allow footnotes in longtable head/foot
\IfFileExists{footnotehyper.sty}{\usepackage{footnotehyper}}{\usepackage{footnote}}
\makesavenoteenv{longtable}
\usepackage{graphicx}
\makeatletter
\newsavebox\pandoc@box
\newcommand*\pandocbounded[1]{% scales image to fit in text height/width
  \sbox\pandoc@box{#1}%
  \Gscale@div\@tempa{\textheight}{\dimexpr\ht\pandoc@box+\dp\pandoc@box\relax}%
  \Gscale@div\@tempb{\linewidth}{\wd\pandoc@box}%
  \ifdim\@tempb\p@<\@tempa\p@\let\@tempa\@tempb\fi% select the smaller of both
  \ifdim\@tempa\p@<\p@\scalebox{\@tempa}{\usebox\pandoc@box}%
  \else\usebox{\pandoc@box}%
  \fi%
}
% Set default figure placement to htbp
\def\fps@figure{htbp}
\makeatother


% definitions for citeproc citations
\NewDocumentCommand\citeproctext{}{}
\NewDocumentCommand\citeproc{mm}{%
  \begingroup\def\citeproctext{#2}\cite{#1}\endgroup}
\makeatletter
 % allow citations to break across lines
 \let\@cite@ofmt\@firstofone
 % avoid brackets around text for \cite:
 \def\@biblabel#1{}
 \def\@cite#1#2{{#1\if@tempswa , #2\fi}}
\makeatother
\newlength{\cslhangindent}
\setlength{\cslhangindent}{1.5em}
\newlength{\csllabelwidth}
\setlength{\csllabelwidth}{3em}
\newenvironment{CSLReferences}[2] % #1 hanging-indent, #2 entry-spacing
 {\begin{list}{}{%
  \setlength{\itemindent}{0pt}
  \setlength{\leftmargin}{0pt}
  \setlength{\parsep}{0pt}
  % turn on hanging indent if param 1 is 1
  \ifodd #1
   \setlength{\leftmargin}{\cslhangindent}
   \setlength{\itemindent}{-1\cslhangindent}
  \fi
  % set entry spacing
  \setlength{\itemsep}{#2\baselineskip}}}
 {\end{list}}
\usepackage{calc}
\newcommand{\CSLBlock}[1]{\hfill\break\parbox[t]{\linewidth}{\strut\ignorespaces#1\strut}}
\newcommand{\CSLLeftMargin}[1]{\parbox[t]{\csllabelwidth}{\strut#1\strut}}
\newcommand{\CSLRightInline}[1]{\parbox[t]{\linewidth - \csllabelwidth}{\strut#1\strut}}
\newcommand{\CSLIndent}[1]{\hspace{\cslhangindent}#1}



\setlength{\emergencystretch}{3em} % prevent overfull lines

\providecommand{\tightlist}{%
  \setlength{\itemsep}{0pt}\setlength{\parskip}{0pt}}



 


\usepackage{mathtools}
\KOMAoption{captions}{tableheading}
\makeatletter
\@ifpackageloaded{tcolorbox}{}{\usepackage[skins,breakable]{tcolorbox}}
\@ifpackageloaded{fontawesome5}{}{\usepackage{fontawesome5}}
\definecolor{quarto-callout-color}{HTML}{909090}
\definecolor{quarto-callout-note-color}{HTML}{0758E5}
\definecolor{quarto-callout-important-color}{HTML}{CC1914}
\definecolor{quarto-callout-warning-color}{HTML}{EB9113}
\definecolor{quarto-callout-tip-color}{HTML}{00A047}
\definecolor{quarto-callout-caution-color}{HTML}{FC5300}
\definecolor{quarto-callout-color-frame}{HTML}{acacac}
\definecolor{quarto-callout-note-color-frame}{HTML}{4582ec}
\definecolor{quarto-callout-important-color-frame}{HTML}{d9534f}
\definecolor{quarto-callout-warning-color-frame}{HTML}{f0ad4e}
\definecolor{quarto-callout-tip-color-frame}{HTML}{02b875}
\definecolor{quarto-callout-caution-color-frame}{HTML}{fd7e14}
\makeatother
\makeatletter
\@ifpackageloaded{caption}{}{\usepackage{caption}}
\AtBeginDocument{%
\ifdefined\contentsname
  \renewcommand*\contentsname{Table of contents}
\else
  \newcommand\contentsname{Table of contents}
\fi
\ifdefined\listfigurename
  \renewcommand*\listfigurename{List of Figures}
\else
  \newcommand\listfigurename{List of Figures}
\fi
\ifdefined\listtablename
  \renewcommand*\listtablename{List of Tables}
\else
  \newcommand\listtablename{List of Tables}
\fi
\ifdefined\figurename
  \renewcommand*\figurename{Figure}
\else
  \newcommand\figurename{Figure}
\fi
\ifdefined\tablename
  \renewcommand*\tablename{Table}
\else
  \newcommand\tablename{Table}
\fi
}
\@ifpackageloaded{float}{}{\usepackage{float}}
\floatstyle{ruled}
\@ifundefined{c@chapter}{\newfloat{codelisting}{h}{lop}}{\newfloat{codelisting}{h}{lop}[chapter]}
\floatname{codelisting}{Listing}
\newcommand*\listoflistings{\listof{codelisting}{List of Listings}}
\makeatother
\makeatletter
\makeatother
\makeatletter
\@ifpackageloaded{caption}{}{\usepackage{caption}}
\@ifpackageloaded{subcaption}{}{\usepackage{subcaption}}
\makeatother
\usepackage{bookmark}
\IfFileExists{xurl.sty}{\usepackage{xurl}}{} % add URL line breaks if available
\urlstyle{same}
% fallback for those not using the hyperref driver hyperxmp:
\makeatletter
\@ifundefined{xmpquote}{\newcommand{\xmpquote}[1]{#1}}{}
\makeatother
\hypersetup{
  pdftitle={SIFT: Search Conjoint},
  pdfauthor={Dan Yavorsky; Eric T. Bradlow},
  pdfkeywords={\xmpquote{consumer search, sequential search, search
costs, conjoint analysis, preference measurement, consideration sets,
choice-based conjoint, structural estimation}},
  colorlinks=true,
  linkcolor={blue},
  filecolor={Maroon},
  citecolor={Blue},
  urlcolor={Blue},
  pdfcreator={LaTeX via pandoc}}


\title{SIFT: Search Conjoint}
\usepackage{etoolbox}
\makeatletter
\providecommand{\subtitle}[1]{% add subtitle to \maketitle
  \apptocmd{\@title}{\par {\large #1 \par}}{}{}
}
\makeatother
\subtitle{Measuring Preferences and Search Costs in a Single Survey
Instrument}
\author{Dan Yavorsky \and Eric T. Bradlow}
\date{2026-08-26}
\begin{document}
\maketitle
\begin{abstract}
Conjoint analysis measures what consumers prefer among alternatives they
have been shown. It is silent on what they would have looked at, what
looking costs them, and what they never see at all. We introduce SIFT, a
survey instrument in which the respondent faces a list-screen of
alternatives, clicks through to inspect individual listings at will, and
then chooses --- a task that reproduces the two-stage list-then-detail
architecture of the environments in which many category purchases are
actually made. Each respondent completes a sequence of such tasks, and
position, list-screen content, and detail-page content are randomized by
design in every one. We fit the resulting click-and-choose data with a
sequential search model in the tradition of Weitzman (1979), recovering
both preference parameters, as conjoint does, and search costs, for
which conjoint has no analogue. Two features distinguish the instrument
from the structural search literature it borrows from. First, the
within-respondent panel of search spells supports individual-level
estimates of preferences and search costs jointly, where field
applications typically observe a single spell per consumer and identify
heterogeneity cross-sectionally. Second, the exogenous variation that
field studies obtain only rarely --- randomized position and information
architecture --- is designed in rather than found. The framework nests
conjoint: as the search cost goes to zero every alternative is
inspected, choice is made over fully realized utilities, and the model
collapses to a standard choice-based conjoint model. Whether the search
cost is distinguishable from zero is therefore a testable question
rather than a maintained assumption, and the answer tells the analyst
whether consideration or preference is the binding constraint in a
category. We derive the likelihood, characterize what the designed
variation buys for identification, and provide aggregate and
hierarchical Bayesian estimation routines with replication code in R. A
simulation study documents parameter recovery at commercial task counts,
and an empirical application fielded in partnership with a consumer
insights consultancy compares SIFT's preference estimates against
conjoint part-worths from a split sample, tests whether estimated search
costs are distinguishable from zero, and decomposes non-purchase into
never-inspected and inspected-but-rejected --- a distinction conjoint
cannot draw, and one that carries different managerial remedies.
\end{abstract}


  \newcommand{\iid}{\stackrel{iid}{\sim}}
  \newcommand{\bfbeta}{\boldsymbol{\beta}}
  \newcommand{\bfgamma}{\boldsymbol{\gamma}}
  \newcommand{\bfkappa}{\boldsymbol{\kappa}}
  \newcommand{\bfeta}{\boldsymbol{\eta}}
  \newcommand{\bftheta}{\boldsymbol{\theta}}
  \newcommand{\bfOmega}{\boldsymbol{\Omega}}
  \newcommand{\bfX}{\mathbf{X}}
  \newcommand{\bfL}{\mathbf{L}}
  \newcommand{\bfZ}{\mathbf{Z}}
  \newcommand{\Jset}{\mathcal{J}}
  \newcommand{\Sset}{\mathcal{S}}
  \newcommand{\Sbar}{\bar{\mathcal{S}}}
  \newcommand{\resv}{z}
  \newcommand{\cost}{c}

\section{Introduction}\label{sec-intro}

\section{Related Literature}\label{sec-lit}

\begin{tcolorbox}[enhanced jigsaw, arc=.35mm, bottomrule=.15mm, bottomtitle=1mm, breakable, colback=white, colbacktitle=quarto-callout-warning-color!10!white, colframe=quarto-callout-warning-color-frame, coltitle=black, left=2mm, leftrule=.75mm, opacityback=0, opacitybacktitle=0.6, rightrule=.15mm, title=\textcolor{quarto-callout-warning-color}{\faExclamationTriangle}\hspace{0.5em}{Citation verification status --- remove before circulation}, titlerule=0mm, toprule=.15mm, toptitle=1mm]

Every reference in this section was reconstructed from memory without
database access and is \textbf{unverified}. Confidence flags follow each
entry: \textbf{{[}H{]}} confident the paper exists and the details are
roughly right; \textbf{{[}M{]}} confident the paper exists, unsure of
journal, year, or coauthors; \textbf{{[}L{]}} believe something like
this exists, may be conflating sources. Treat \textbf{{[}L{]}} items as
search leads, not references. Nothing here should reach a referee, a
proposal, or a client before it is checked against the record. The three
entries marked \textbf{novelty gate} determine whether the paper's
central claim survives; read those first.

\end{tcolorbox}

\subsection{Optimal Sequential Search}\label{sec-lit-theory}

The theoretical backbone. The instrument's data-generating process is
the Weitzman protocol run inside a survey.

\begin{itemize}
\tightlist
\item
  \textbf{Weitzman (1979), ``Optimal Search for the Best Alternative,''
  \emph{Econometrica}. {[}H{]}} The Pandora's box problem, and the model
  the instrument is built to estimate. Each box has a known prize
  distribution and a known opening cost. The optimal policy factors into
  three parts: a \emph{reservation value} computed for each box
  independently of every other box; a \emph{selection rule} (open boxes
  in descending reservation value); and a \emph{stopping rule} (stop
  when the best realized prize exceeds the highest remaining reservation
  value). The independence of the reservation values is what makes the
  model estimable at survey scale.
\item
  \textbf{Stigler (1961), ``The Economics of Information,''
  \emph{Journal of Political Economy}. {[}H{]}} The fixed-sample-size
  alternative: the consumer commits to \(n\) searches up front rather
  than deciding adaptively. The competing protocol, and the one SIFT's
  panel structure may be able to test against sequential search
  respondent by respondent.
\item
  \textbf{McCall (1970), \emph{Quarterly Journal of Economics}. {[}M{]}}
  Reservation-wage formulation of sequential search with recall; the
  labor-search ancestor of the marketing applications below.
\end{itemize}

\emph{Simultaneous versus sequential.} These are two competing
protocols, not two descriptions of one thing. Sequential search fits
click-through-and-inspect; fixed-sample search fits ``get three
quotes.'' Honka and Chintagunta (2017) is the reference on
distinguishing them empirically, and is directly relevant here because
SIFT's within-respondent panel could run that test at the individual
level rather than the population level.

\subsection{Structural Search with Field and Clickstream
Data}\label{sec-lit-field}

Context rather than precedent: all of this work is observational or
platform data with no designed instrument. SIFT's credibility with an
economics audience depends on connecting to it cleanly.

\begin{itemize}
\tightlist
\item
  \textbf{Hong and Shum (2006), \emph{RAND Journal of Economics}.
  {[}H{]}} Recovers search costs from price distributions alone, without
  search data. The extreme case of how little the field has typically
  had to work with.
\item
  \textbf{Hortacsu and Syverson (2004), \emph{Quarterly Journal of
  Economics}. {[}H{]}} S\&P 500 index funds. Search frictions with
  differentiated products, in a setting where the products are nearly
  identical and the frictions are not.
\item
  \textbf{Kim et al. (2010), \emph{Marketing Science}. {[}H{]}} Amazon
  view-rank data for camcorders. One of the first structural search
  models fit to online browsing data in marketing.
\item
  \textbf{De los Santos et al. (2012), ``Testing Models of Consumer
  Search Using Data on Web Browsing,'' \emph{American Economic Review}.
  {[}H{]}} Comscore browsing data; tests sequential search against
  fixed-sample search and finds fixed-sample fits better in their
  setting. The empirical warning shot for assuming the Weitzman protocol
  rather than testing it.
\item
  \textbf{Koulayev (2014), \emph{RAND Journal of Economics}. {[}M{]}}
  Hotel search with click data; identification and estimation for
  differentiated products.
\item
  \textbf{Honka (2014), \emph{RAND Journal of Economics}. {[}H{]}} Auto
  insurance; separates search costs from switching costs. The template
  for arguing that an estimated friction is the friction you claim it
  is.
\item
  \textbf{Honka and Chintagunta (2017), \emph{Marketing Science}.
  {[}M{]}} Simultaneous or sequential; see above.
\item
  \textbf{Chen and Yao (2017), \emph{Management Science}. {[}M{]}}
  Sequential search with refinement --- filters and sorts modeled as
  search actions --- on Expedia clickstream data. Closest existing work
  in spirit to a list-screen interface, and the natural precedent if the
  instrument later adds filtering.
\item
  \textbf{Ursu (2018), ``The Power of Rankings,'' \emph{Marketing
  Science}. {[}H{]}} Expedia's randomized-position experiment. The
  cleanest source of exogenous variation in the field literature,
  obtained once, as a rare gift from a platform. The contrast with SIFT
  is direct: here the same variation is free, designed in, and present
  in every task.
\item
  \textbf{Ursu et al. (2020), \emph{Marketing Science}. {[}M{]}} Search
  duration and time spent, rather than search incidence alone.
\item
  \textbf{Yavorsky et al. (2021), \emph{Quantitative Marketing and
  Economics}. {[}H{]}} Dealership visits in the U.S. auto market;
  identification of search costs from observed search sets. (Author's
  own; the estimation code is an implementation asset, see
  Section~\ref{sec-estimation}.)
\item
  \textbf{Bronnenberg et al. (2016), ``Zooming In on Choice,''
  \emph{Journal of Marketing Research}. {[}M{]}} Descriptive clickstream
  evidence on how consumers actually move through product pages.
  Face-validity evidence for the interface design in

  \begin{enumerate}
  \def\labelenumi{\arabic{enumi}.}
  \tightlist
  \item
  \end{enumerate}
\item
  \textbf{2, \emph{Quantitative Marketing and Economics}. {[}M{]} ---
  novelty gate} Framework and consolidation review of the sequential
  search model. Best single entry point to the current state of model
  specification, identification, and estimation practice; read first to
  establish what is already settled.
\item
  \textbf{Greminger (2022), \emph{Management Science}. {[}M{]}} Search
  with discovery of alternatives --- the consumer does not know the full
  choice set at the outset. Relevant to how much of the list-screen the
  instrument reveals up front.
\end{itemize}

\subsection{Lab and Experimental Evidence on Search}\label{sec-lit-lab}

Two experimental literatures that have barely spoken to each other, and
the direct methodological precedent for running search as a designed
task.

\subsubsection{Does anyone follow the optimal
rule?}\label{does-anyone-follow-the-optimal-rule}

\begin{itemize}
\tightlist
\item
  \textbf{SCHOTTER and BRAUNSTEIN (1981), \emph{Economic Inquiry}.
  {[}M{]}} Early sequential-search experiments with real incentives.
\item
  \textbf{Hey (1987), ``Still Searching,'' \emph{Journal of Economic
  Behavior and Organization}. {[}M{]}} Subjects use rules of thumb
  rather than optimal reservation strategies.
\item
  \textbf{Cox and Oaxaca (1989). {[}M{]}} Finite-horizon search; broadly
  supportive of reservation-rule behavior with noise. The more favorable
  read.
\item
  \textbf{Sonnemans (1998). {[}M{]}} Verbal-protocol evidence and
  classification of search strategies.
\item
  \textbf{Brown et al. (2011). {[}L{]}} Real-time search in the
  laboratory and in the market.
\end{itemize}

\emph{Net read.} Mixed, and this is the single largest scientific risk
to the project. Behavior is directionally consistent with
reservation-value logic, but the deviations are systematic rather than
random: too little search relative to the optimum, order effects, and
sensitivity to sunk costs. The paper should state this openly and early
rather than let a referee find it.

\subsubsection{Information acquisition: Mouselab and
eye-tracking}\label{information-acquisition-mouselab-and-eye-tracking}

Mechanically, the SIFT task is a Mouselab information board wearing a
commercial skin.

\begin{itemize}
\tightlist
\item
  \textbf{Payne et al. (1988); Payne et al. (1993), \emph{The Adaptive
  Decision Maker}. {[}H{]}} The Mouselab paradigm: information boards in
  which subjects click cells to reveal attribute values, with the
  acquisition sequence recorded. Descriptive and process-focused rather
  than structural.
\item
  \textbf{Gabaix et al. (2006), ``Costly Information Acquisition,''
  \emph{American Economic Review}. {[}H{]} --- novelty gate} A
  Mouselab-style experiment paired with a \emph{structural}
  boundedly-rational model (directed cognition) estimated on the
  acquisition data. The closest existing template for the core maneuver
  --- run an information-acquisition experiment, fit a structural
  search-type model to the acquisition record. Not Weitzman, but the
  same move.
\item
  \textbf{Caplin et al. (2011), ``Search and Satisficing,''
  \emph{American Economic Review}. {[}H{]}} Choice-process data showing
  that people stop before finding the best available option. The
  satisficing alternative against which model fit should be reported.
\item
  \textbf{Reutskaja et al. (2011), \emph{American Economic Review}.
  {[}H{]}} Eye-tracking evidence on search under time pressure with
  supermarket-style displays.
\end{itemize}

\emph{Net read.} The experimental apparatus is well validated and has
been for decades. What has largely not been done is fitting a Weitzman
model to acquisition data and reporting search costs in interpretable
units.

\subsection{Information Search inside Preference
Measurement}\label{sec-lit-preference}

The closest existing work in marketing, and the location of the sharpest
threat to novelty.

\begin{itemize}
\tightlist
\item
  \textbf{Yang et al. (2015), ``A Bounded Rationality Model of
  Information Search and Choice in Preference Measurement,''
  \emph{Journal of Marketing Research}. {[}M{]} --- novelty gate} Models
  within-task information search during a conjoint-style task using
  eye-tracking data, with an optimal-stopping flavor. If one paper
  threatens the contribution, it is this one. Read it first and read it
  carefully.
\item
  \textbf{Yang et al. (2018), \emph{Marketing Science}. {[}M{]}}
  Attention and information processing in incentive-aligned choice
  experiments.
\item
  \textbf{Meißner et al. (2016), \emph{Journal of Marketing Research}.
  {[}M{]}} Eye-tracking in conjoint; processing becomes measurably more
  efficient with practice. Directly relevant because a multi-task panel
  will exhibit learning across tasks, which appears in the model as
  drift in the search-cost parameter (Section~\ref{sec-lit-lab},
  Section~\ref{sec-discussion}).
\end{itemize}

\emph{Distinction to preserve --- and to verify before relying on it.}
These papers study search \emph{within} a task, over the attributes of
profiles already displayed, mostly in service of better preference
estimates. SIFT studies search \emph{across alternatives}, in a
two-stage list-then-detail architecture, where the search cost is itself
a deliverable rather than a nuisance parameter. Related, but not the
same object. Confirm this holds before building on it.

\subsection{Consideration Sets and Screening Rules in
Conjoint}\label{sec-lit-consideration}

The existing marketing answer to ``people do not evaluate everything,''
and the comparison a conjoint-literate reader will reach for first.

\begin{itemize}
\tightlist
\item
  \textbf{Gilbride and Allenby (2004), \emph{Marketing Science}.
  {[}H{]}} A choice model with conjunctive, disjunctive, and
  compensatory screening rules. The incumbent formulation.
\item
  \textbf{Hauser et al. (2010), \emph{Journal of Marketing Research}.
  {[}H{]}} Disjunctions of conjunctions; cognitive simplicity in
  consideration rules.
\item
  \textbf{Yee et al. (2007), \emph{Marketing Science}. {[}M{]}}
  Greedoid-based noncompensatory inference.
\item
  \textbf{Dzyabura and Hauser (2011), \emph{Marketing Science}. {[}M{]}}
  Active machine learning for consideration heuristics.
\item
  \textbf{Hauser (2014), \emph{Journal of Business Research}. {[}M{]}}
  Review of consideration-set heuristics.
\end{itemize}

\emph{Sharpest one-line contrast.} This literature \emph{infers}
screening rules from choice outcomes, treating consideration as latent
structure. SIFT \emph{observes} the consideration process directly and
prices it in units of search cost. Inference versus observation.

\subsection{Conjoint Practice, Incentive Alignment, and Commercial
Precedent}\label{sec-lit-practice}

\begin{itemize}
\tightlist
\item
  \textbf{Allenby et al. (2019), \emph{Handbook of the Economics of
  Marketing}. {[}H{]}} Economic foundations of conjoint analysis; the
  statement of the workhorse model that SIFT nests.
\item
  \textbf{Marshall and Bradlow (2002), \emph{JASA}. {[}H{]}} A unified
  approach to conjoint analysis models; precedent for treating distinct
  response formats as one likelihood family.
\item
  \textbf{Ding et al. (2005), \emph{Journal of Marketing Research}; Ding
  (2007), \emph{Journal of Marketing Research}. {[}H{]}} Incentive
  alignment in conjoint. If estimated search costs are to carry units
  that mean anything, the task needs real stakes; this is the precedent
  for how to supply them.
\item
  \textbf{Sawtooth ACBC (adaptive choice-based conjoint).} Not an
  academic reference, but the closest commercial product to a two-stage
  consideration instrument: a screening section (must-haves,
  unacceptables) followed by a choice tournament. Study as design
  precedent, and as the thing a client may already have bought.
\item
  \textbf{Virtual-shelf and shopper-simulation vendors.} Several field
  shelf-set exercises with click data, reported descriptively (time to
  first click, click share). No known case of backing a structural
  search model out of one. \textbf{This is where prior art most likely
  hides; check directly.}
\end{itemize}

\subsection{Positioning}\label{sec-lit-positioning}

\section{The SIFT Instrument}\label{sec-instrument}

\subsection{Task Architecture}\label{sec-instrument-architecture}

A SIFT task presents the respondent with two kinds of screen, arranged
the way online retail arranges them. The first is a \textbf{listing
page}: a set of \(J\) alternatives displayed as an ordered vertical
list, each occupying one row. The second is a \textbf{product page}, one
per alternative, reached by clicking its row. The pair corresponds to
what commercial practice calls the product listing page and the product
detail page, and the correspondence is deliberate --- the instrument is
meant to look to the respondent like the environments in which the
category is actually shopped.

The two screens differ in what they show. Every alternative is described
by the same set of attributes, but the attributes are partitioned. A
small number are displayed on the listing page for all \(J\)
alternatives at once --- position in the list, brand, and price are the
natural candidates, and are the ones that list pages typically carry in
practice. The remainder are shown only on an alternative's product page.
A respondent who reads the listing page and clicks nothing therefore
knows a little about every alternative; a respondent who opens a product
page knows everything about that one alternative and still only a little
about the rest.

This partition is what gives the task its structure. Listing-page
attributes are free: the respondent acquires them by looking at the
screen she is already on. Product-page attributes are costly: acquiring
them requires a click, a page load, and the time to read what the page
contains. The respondent chooses how many of those costs to pay and in
what order.

The protocol within a task is as follows. The respondent begins on the
listing page. She may open any alternative's product page, return to the
listing page, and open another, as many times as she likes and in any
order she likes. Every alternative she has already opened remains
available to her --- she may return to a previously viewed product page
at no further cost, and she may select an alternative she viewed several
clicks ago. The task ends in one of two ways. She selects an alternative
to purchase, which she may do only from a product page she has opened,
so that no alternative can be purchased sight-unseen. Or she declines to
purchase anything and exits the task, which she may do at any point,
including immediately, without opening a single product page.

Two features of this protocol deserve emphasis because the model in
Section~\ref{sec-model} rests on them. First, selection requires prior
inspection, which is both realistic --- one buys from the product page,
not from the list --- and what allows the observed choice to be read as
a choice over resolved rather than expected utilities. Second, recall is
free and unlimited, so a respondent is never punished for having looked
at something and moved on. Together these place the task squarely in the
Weitzman (1979) protocol: the listing page is a set of closed boxes
whose contents are unknown but whose distributions are not, a click is
the act of opening one at a cost, and the decision to stop and select is
the decision to take the best prize already in hand rather than pay to
open another. The no-purchase exit is the outside option, which the
respondent may take whether or not she has searched.

The exit option matters for what the instrument can measure. A
respondent who searches nothing and leaves, a respondent who searches
half the list and leaves, and a respondent who searches half the list
and buys are producing three different pieces of evidence about the same
two parameters. Conjoint's no-choice option records only the first
distinction between them.

Respondents complete a sequence of tasks rather than one, exactly as in
choice-based conjoint. Each task is a fresh listing page with a fresh
set of alternatives, and the respondent searches and chooses within it
independently. The current design fields ten tasks per respondent. This
is the feature that separates the instrument most sharply from the field
literature it borrows from. A platform dataset typically observes one
search spell per consumer, so preferences and search costs must be told
apart across consumers. Here each respondent supplies a panel of search
spells, and the same respondent's behavior across tasks can be used to
separate what she likes from what looking costs her.

\subsection{Randomization and Experimental
Design}\label{sec-instrument-randomization}

Three things vary from task to task, and all three vary by design rather
than by circumstance.

\textbf{Position.} The order in which the \(J\) alternatives appear on
the listing page is randomized independently of everything else about
them. An alternative that appears first in one task appears fourth in
another, with the same attributes and the same competitors. Position is
thus uncorrelated by construction with quality, price, brand, and every
other attribute --- the correlation that makes observational ranking
data so difficult to use, since platforms rank on relevance and
relevance is a function of exactly the things that drive utility. Field
work has obtained this variation only rarely and only as a gift, most
cleanly in Ursu (2018). Here it is free and present in every task.

\textbf{Composition.} The set of alternatives available on the listing
page varies across tasks: which alternatives appear, and how many. This
is the analogue of varying the choice set in a conjoint design, and it
serves the same purpose of placing each alternative against varied
competition, but it does additional work here because the attractiveness
of the alternatives a respondent has \emph{not} yet opened is what
determines whether she keeps searching.

\textbf{Attribute levels.} The attributes themselves are fixed across
tasks --- the same characteristics describe every alternative in every
task, as in conjoint --- while the levels those attributes take on are
varied across alternatives and across tasks according to an experimental
design. This is conventional conjoint practice and needs no defense;
what is new is only that some of these attributes are revealed on the
listing page and the rest behind a click.

A fourth lever is available and is not yet part of the design: the
\emph{partition} itself --- which attributes appear on the listing page
and which are held back for the product page --- could be randomized
across respondents rather than fixed. Doing so would make the
information architecture an experimental treatment, and would let the
instrument speak to what belongs on a list page versus a detail page. It
also complicates identification, since it changes what the respondent
knows before searching. We flag it here as a design option and return to
it in Section~\ref{sec-managerial}.

We defer the question of what each of these three sources of variation
identifies to Section~\ref{sec-estimation}, and the question of how they
enter the likelihood to Section~\ref{sec-model}. The point to carry
forward from this section is only that the variation is designed rather
than found: it is orthogonal by construction, it is present in every
task rather than in a single natural experiment, and it costs nothing to
produce.

\subsection{Data Recorded}\label{sec-instrument-data}

The instrument records the full interaction, not merely its outcome. For
each respondent \(i\) and task \(t\) we retain three layers.

\textbf{The design.} The complete attribute matrix for all \(J\)
alternatives shown, whether or not the respondent ever saw them; the
position each alternative occupied on the listing page; and the
partition of attributes into listing-page and product-page sets. The
design is known without error because it was generated, which is a
meaningful advantage over field data, where the analyst typically knows
what was clicked but must reconstruct what was displayed.

\textbf{The event stream.} Every action the respondent takes, in order,
with a timestamp: each product-page opening (which alternative, at which
position, at what time), each return to the listing page, each
re-opening of a previously viewed product page, and the terminal event
--- a purchase of a named alternative or an exit without purchase. The
event stream is the raw object; everything below is derived from it.

\textbf{Derived search quantities.} From the event stream we construct
the objects the model in Section~\ref{sec-model} takes as data: the
\emph{search set}, meaning which alternatives were opened; the
\emph{search order}, meaning the sequence in which they were opened; the
\emph{stopping point}, meaning how many alternatives were opened before
the respondent stopped; and the \emph{choice}, meaning which member of
the search set was purchased, or that none was. Timing data yield a
parallel set of duration measures --- time to first click, dwell time
per product page, total task duration --- which are informative about
search effort in the sense of Ursu et al. (2020) and which will also be
needed to net out any imposed latency from voluntary reading time.

It is worth stating plainly what this means for the alternatives a
respondent never opened. They are not missing data. Under the model they
are the outcome of a decision --- the respondent judged that opening
them was not worth the cost, given what she had already found --- and
they carry information about both her preferences and her search cost.
Non-inspection is the observation that conjoint has no way to record,
and recovering it is much of the point of the instrument.

\section{Model}\label{sec-model}

\subsection{Setup and Notation}\label{sec-model-setup}

We follow the notation of 2 wherever the two are compatible, so that
readers arriving from the sequential-search literature can map our
objects onto theirs without translation. Two departures are forced on us
by the instrument, and one by the object being searched over; we flag
all three below. Section~\ref{sec-appendix-notation} gives a crosswalk
to 2 and Yavorsky et al. (2021) for readers coming from either.

\textbf{Indices.} Respondent \(i = 1,\dots,N\) completes tasks
\(t = 1,\dots,T\). Task \(t\) presents respondent \(i\) with a listing
page holding the alternatives \(\Jset_{it} = \{1,\dots,J_{it}\}\),
together with the no-purchase option \(j=0\), which is always available
and never requires a click. The index \(j\) names an \emph{alternative},
not its rank and not its screen position: because position is randomized
(Section~\ref{sec-instrument-randomization}), the paper must hold three
things apart at once --- an alternative's identity \(j\), the position
it occupied on the listing page in task \(t\), and its rank in the
respondent's reservation-value ordering. Search order is carried by a
separate index \(h\), where \(h=1\) denotes the first product page the
respondent opened.

\textbf{Utility.} Respondent \(i\)'s utility from alternative \(j\) in
task \(t\) is

\begin{equation}\protect\phantomsection\label{eq-utility}{
u_{itj} \;=\; \underbrace{\bfX_{itj}'\bfbeta_i + \mu_{itj}}_{\displaystyle \delta_{itj}\ \text{(known before clicking)}} \;+\; \underbrace{\bfL_{itj}'\bfkappa_i + \varepsilon_{itj}}_{\displaystyle \text{(revealed by clicking)}} ,
}\end{equation}

where \(\bfX_{itj}\) collects the attributes displayed on the listing
page, \(\bfL_{itj}\) the attributes displayed only on the product page,
\(\mu_{itj}\) is a pre-search taste shock observed by the respondent but
not by us, and \(\varepsilon_{itj}\) is a post-search taste shock. Write
\(\xi_{itj} = \bfX_{itj}'\bfbeta_i\) for the observed part of pre-search
utility and \(\delta_{itj} = \xi_{itj} + \mu_{itj}\) for pre-search
utility. Both shocks are mean-zero normal,
\(\mu_{itj} \iid N(0,\sigma_\mu^2)\) and
\(\varepsilon_{itj} \iid N(0,\sigma_\varepsilon^2)\), with
\(\sigma_\mu = 1\) imposed as the scale normalization. The no-purchase
option has \(u_{it0} = \bfX_{it0}'\bfbeta_i + \mu_{it0}\) with no
post-search component: nothing about it is learned by clicking, so its
utility is resolved from the outset.

The term \(\bfL_{itj}'\bfkappa_i\) is the departure that matters most.
In the match-value formulations that dominate empirical work, everything
learned through search is a scalar shock \(\varepsilon_{itj}\)
unobserved by the researcher both before and after search. Here the
product page reveals \emph{attributes we designed}, so part of what
search delivers is observed, priced by \(\bfkappa_i\), and ---
critically --- drawn from a distribution the analyst chose rather than
assumed. 2 give exactly this general form in their §2.4 and note it is
uncommon in practice; SIFT is a case where it is unavoidable, since
recovering \(\bfkappa_i\) for product-page attributes is half of what
the instrument is for. The reservation-value discussion below takes up
what this costs computationally.

\textbf{Search costs.} Opening alternative \(j\)'s product page costs

\begin{equation}\protect\phantomsection\label{eq-cost}{
\cost_{itj} \;=\; \exp\!\left(\bfZ_{itj}'\bfgamma_i\right),
}\end{equation}

exponentiated to enforce positivity, as is standard. \(\bfZ_{itj}\)
holds a respondent-level intercept and the randomized listing-page
position, and may hold listing-page attributes that plausibly shift
salience. Overlap between \(\bfX\) and \(\bfZ\) is permitted;
Section~\ref{sec-estimation} takes up what separates the two.

\textbf{Where the randomness lives.} Reservation values must carry a
stochastic component, or search order would be a deterministic function
of observables and no observed order could be rationalized (Ursu et al.
2024). Two specifications supply one. The first, used above and by most
of the literature, puts a pre-search taste shock \(\mu_{itj}\) in
utility. The second, due to Chung et al. (2024), drops \(\mu_{itj}\) and
instead treats the search cost itself as random at the
respondent-alternative level, with \(\cost_{itj}\) drawn from a
distribution whose location carries the position effect. The two are
observationally similar but econometrically quite different: under the
first, the scale of reservation utilities and the scale of realized
utilities are jointly determined by \(\sigma_\mu\) and
\(\sigma_\varepsilon\), and estimating both requires an auxiliary step;
under the second, search-cost dispersion fixes the scale of reservation
values while \(\sigma_\varepsilon\) alone fixes the scale of realized
utilities, and everything is estimable in one pass.

2 §4.4.1 object to the second on interpretive grounds, asking why the
cost of clicking a product detail page would vary at the
consumer-product level. The objection has least force here, because we
randomize the thing that makes it vary: the same alternative sits at a
different position for different respondents and in different tasks, so
respondent-alternative variation in click cost is designed in rather
than assumed. Section~\ref{sec-estimation} takes up which specification
to estimate; Chung et al. (2024) show that relative preference
parameters survive getting this choice wrong, though search-cost
\emph{levels} do not.

\textbf{Search and choice.} Let \(\Sset_{it} \subseteq \Jset_{it}\)
denote the set of alternatives whose product pages respondent \(i\)
opened in task \(t\), \(\Sbar_{it}\) its complement,
\(H_{it} = |\Sset_{it}|\) the number of clicks, and
\(y_{it} \in \Sset_{it} \cup \{0\}\) the alternative purchased, with
\(y_{it}=0\) recording exit without purchase. Recall is free, so
\(y_{it}\) may be any member of \(\Sset_{it}\), and selection requires
prior inspection, so it may not be drawn from \(\Sbar_{it}\).

\textbf{Reservation values.} The reservation value \(\resv_{itj}\)
equates the marginal cost of opening a product page to the marginal
expected benefit,

\begin{equation}\protect\phantomsection\label{eq-reservation}{
\int_{\resv_{itj}}^{\infty} \left(u - \resv_{itj}\right) dF_{itj}(u) \;=\; \cost_{itj},
}\end{equation}

where \(F_{itj}\) is the respondent's belief about \(u_{itj}\) prior to
clicking. Additive separability gives
\(\resv_{itj} = \delta_{itj} + g(\cost_{itj})\), with \(g\) decreasing
in search cost and depending on the belief distribution only through the
post-search component.

Here Equation~\ref{eq-utility} exacts a price, though a smaller one than
it first appears. When the post-search component is a single normal
shock, \(g\) has the familiar closed form
\(g(\cost) = m(\cost/\sigma_\varepsilon)\,\sigma_\varepsilon\) with
\(m\) solving \(\cost/\sigma_\varepsilon = \phi(m) + m[\Phi(m)-1]\) (Kim
et al. 2010), and this is what every existing implementation computes.
With designed product-page attributes the pre-click belief about
\(\bfL_{itj}'\bfkappa_i + \varepsilon_{itj}\) is instead a
\emph{mixture} of normals --- one component per attribute-level
combination the design can produce, with weights fixed by the design ---
and \(g\) has no scalar closed form.

Two properties of a designed instrument keep this cheap. First, the
mixture is \emph{common across alternatives}: a balanced design draws
\(\bfL_{itj}\) from the same distribution at every position, so the
respondent's pre-click belief is identical for every alternative she has
not yet opened, and the mixture enters the model only through \(g\)
rather than alternative by alternative. The same balance delivers the
independence across \(j\) that Weitzman's rules require. Second, the
left side of Equation~\ref{eq-reservation} is then a
probability-weighted sum of normal partial expectations, strictly
decreasing in \(\resv_{itj}\), so one root-find recovers \(g\) for each
distinct search cost --- and because position is categorical, the number
of distinct search costs per respondent is the number of positions, not
the number of alternative-task pairs. The scalar-normal case is the
one-component special case, and nothing in Weitzman's rules requires
normality, so the decision rules of Section~\ref{sec-model-policy} carry
over unchanged.

A variance-matched normal approximation is also available, and it is the
specification we estimate. Absorbing the design mean of
\(\bfL'\bfkappa_i\) into \(\delta_{itj}\) and replacing the mixture with
a normal of variance
\(\tilde\sigma_i^2 = \bfkappa_i'\mathrm{Var}(\bfL)\bfkappa_i + \sigma_\varepsilon^2\)
returns the model exactly to the Kim closed form with an inflated
post-search standard deviation. This is more than a computational
convenience. The benefit of search --- the object the literature
normalizes away for want of variation to estimate it (Ursu 2018; Morozov
et al. 2021) --- is here \emph{partly known}, because
\(\mathrm{Var}(\bfL)\) is fixed by the design rather than assumed.
Because \(g\) then depends on its arguments only through
\(\cost/\tilde\sigma_i\), a single one-dimensional inverse serves every
respondent, which is what keeps estimation tractable at panel scale
(Section~\ref{sec-est-estimation}). We report the exact mixture as a
generalization and use it as a robustness check.

\textbf{Parameters.} Collect the respondent-level parameters in
\(\bftheta_i = (\bfbeta_i, \bfkappa_i, \bfgamma_i)\) and the population
parameters governing their distribution in \(\bfOmega\). We reserve
\(\bftheta\) for the parameter vector throughout and write
\(\log \sigma_\varepsilon\) explicitly where the log scale is meant.

\subsection{Optimal Search Policy}\label{sec-model-policy}

\subsection{Likelihood}\label{sec-model-likelihood}

\subsection{Conjoint as the Zero-Search-Cost Special
Case}\label{sec-model-nesting}

\subsection{Alternative Search Protocols}\label{sec-model-alternatives}

\section{Identification and Estimation}\label{sec-estimation}

Search models are identified from patterns that ordinary choice data
cannot produce: \emph{which} alternatives a respondent opened, \emph{in
what order}, \emph{how many}, and \emph{which} of the opened
alternatives was ultimately bought. This section sets out what each of
those margins identifies, what SIFT's designed variation adds to
arguments that field studies must make on observational grounds, and
what remains normalized rather than estimated.

Throughout, the object of interest is the pair \((\bfbeta, \bfgamma)\)
--- preferences and search costs --- together with the post-search
standard deviation \(\tilde\sigma\), which most of the literature fixes
to one and which we argue must be estimated.

\subsection{What the Designed Variation Buys}\label{sec-est-variation}

\subsubsection{Position enters search costs, not
utility}\label{position-enters-search-costs-not-utility}

The identifying assumption underlying every position-based search-cost
shifter is that a listing's position affects what it costs to inspect an
alternative, and nothing else. Ursu (2018) §4.3 takes this seriously
rather than asserting it, enumerating the ways position could instead
enter the model --- through pre-search utility, through the post-search
spread, or through combinations --- and ruling them out with
observational tests. Because SIFT randomizes position within every task
and records the full click sequence, we can run those tests as
experimental contrasts rather than conditional comparisons.

Two statistics suffice, and neither requires fitting a model.

\textbf{T1. Among clicked alternatives, does position predict purchase?}
\textbf{T2. Is the first click near the top of the page?}

\begin{tcolorbox}[enhanced jigsaw, arc=.35mm, bottomrule=.15mm, bottomtitle=1mm, breakable, colback=white, colbacktitle=quarto-callout-important-color!10!white, colframe=quarto-callout-important-color-frame, coltitle=black, left=2mm, leftrule=.75mm, opacityback=0, opacitybacktitle=0.6, rightrule=.15mm, title=\textcolor{quarto-callout-important-color}{\faExclamation}\hspace{0.5em}{The sign of T1 is the opposite of the natural intuition}, titlerule=0mm, toprule=.15mm, toptitle=1mm]

It is tempting to reason that if position acts only on costs then it
should have no bearing on what gets bought once someone has looked, so a
flat T1 indicates position-in-cost. Simulation says otherwise, and the
reason is selection. When position raises the search cost, an
alternative in a \emph{bad} position must have unusually high pre-search
utility to be worth opening at all. Conditional on having been clicked,
bad-position alternatives are therefore \emph{better} than good-position
ones, and T1 comes out strongly \textbf{positive}. When position instead
enters pre-search utility, the direct effect and the selection effect
roughly cancel and T1 is approximately \textbf{flat}.

In our simulations (400 respondents × 8 tasks) T1 has slope \(+0.39\)
(\(p \approx 5\times10^{-15}\)) when position acts on costs, and
\(-0.05\) (\(p = 0.06\)) when it acts on utility. Reading a flat T1 as
evidence for position-in-cost would get the answer exactly backwards.

\end{tcolorbox}

T2 separates the remaining case. If position widens the post-search
spread, the reservation value \emph{rises} with position, so the
worst-placed alternatives would be opened first. That is what we see:
mean first-click position 4.96 against a uniform benchmark of 3.5,
versus 2.21 and 2.57 in the cost and utility arms.

The two tests together form a decision rule:

{\def\LTcaptype{none} % do not increment counter
\begin{longtable}[]{@{}lll@{}}
\toprule\noalign{}
T2 & T1 & position acts on \\
\midrule\noalign{}
\endhead
\bottomrule\noalign{}
\endlastfoot
first click near top & strongly positive & \textbf{search cost} \\
first click near top & flat & pre-search utility \\
first click near bottom & --- & post-search spread \\
\end{longtable}
}

We regard this as a design check to be run and reported in every SIFT
study rather than assumed. Details, code, and the full three-arm
simulation are in \href{../comparisons/compare-ursu-position.qmd}{the
position comparison}.

\subsubsection{Separating preferences from search
costs}\label{separating-preferences-from-search-costs}

The sharper problem is not whether position shifts costs but whether a
high search cost can be distinguished from a low taste. Morozov et al.
(2021) gives the cleanest statement of the separation (his §4.3.2):
stopping decisions depend on the \emph{realized} utilities of
alternatives already searched, but not on their search costs, which are
sunk by the time the stopping decision is made. Trade the utility
intercept against the search cost so as to hold the mean reservation
value fixed, and the search \emph{order} is unchanged --- but the higher
intercept raises realized utility and so makes stopping more likely.

His Appendix D sharpens this into an observable moment: the
\textbf{recall rate}, the probability of purchasing an alternative
encountered earlier in the sequence rather than the most recent one,
responds to the utility intercept but not to the search cost. SIFT
observes recall directly, because respondents have free recall and we
record the entire click sequence. This is our cleanest separation
moment, and it is available in the raw data without any modelling.

\subsubsection{What the panel buys, and what it does
not}\label{what-the-panel-buys-and-what-it-does-not}

Morozov et al. (2021) proves nonparametric identification of
individual-level preferences \emph{and} search costs from panel search
data. It is worth being precise about the strength of that result,
because it is easy to overclaim. The argument assumes the number of
sessions per consumer grows without bound; with between 2 and 35
sessions per consumer and a mean of 5, Morozov et al. (2021) estimates
hyperparameters rather than individual parameters (his §5.1).

At the task counts SIFT can realistically field --- call it \(T=10\) ---
we are in the same position, and we say so plainly: what the panel
delivers is \emph{hierarchical posteriors in the conjoint tradition},
not nonparametric individual identification. That is the standard the
conjoint audience already applies to part-worths, and it is the right
standard here.

The panel also does something less obvious. Ursu (2018) reports baseline
search costs above \$102 and attributes the magnitude to observing a
single click per search impression with no way to link a consumer's
impressions; restricting attention to impressions with two or more
clicks cuts mean estimated search costs by up to 83\%. Single-spell data
mechanically inflates search costs. The panel is therefore a
\emph{levels} argument as much as a heterogeneity argument.

\subsubsection{What neither the panel nor the shifter buys
alone}\label{what-neither-the-panel-nor-the-shifter-buys-alone}

The two ingredients are complements, and the literature illustrates this
by having one or the other but never both. Morozov et al. (2021) has the
panel and still normalizes the post-search standard deviation, reporting
that little variation in his data speaks to it (his Appendix C).
Yavorsky et al. (2021) has an exogenous search-cost shifter and no
panel, and does estimate it --- finding \(\tilde\sigma = 8.16\) against
the conventional normalization of one, with a confidence interval that
excludes one and consequences large enough to reverse the sign of one
search-cost coefficient and to understate a counterfactual by a factor
of three to four.

SIFT has both. It also has something neither has: because we design the
product page, part of the post-search component is \emph{fixed by the
design} rather than inferred.

This last point is stronger than it first appears, and it is worth
stating carefully. In a field setting the analyst never observes what
inspection revealed; the post-search shock is a residual, and all that
can be recovered is its variance. In SIFT the analyst \emph{wrote} the
product page. So for every clicked alternative the revealed attributes
\(\bfL_{ijt}\) are known, and they enter realized utility as
\(\bfL_{ijt}'\bfkappa\) --- a mean shift, not a variance. Preferences
over attributes that are only visible \emph{after} a click are therefore
identified in exactly the way conjoint part-worths are: from which
alternative gets bought given what was shown.

The distinction is easy to lose in implementation, with consequences
that are silent. The ex ante spread \(\tilde\sigma\) belongs in the
reservation value, because that is computed before the click; the
realized \(\bfL\) belongs in the utility, because that is evaluated
after it. Use \(\tilde\sigma\) for both and \(\bfkappa\) enters the
likelihood only through
\(\tilde\sigma^2 = \bfkappa'\mathrm{Var}(\bfL)\bfkappa + \sigma_\xi^2\),
where it trades off exactly against \(\sigma_\xi\) and is not identified
at all --- while every other parameter continues to recover normally, so
nothing looks wrong. This is the sharpest version of the identification
claim in the paper and it belongs in the introduction.

\subsection{Normalization}\label{sec-est-normalization}

Setting the standard deviation of the pre-search shock to one fixes the
scale of utility, as in 2 §4.4.2 and Morozov et al. (2021) Appendix C.
That normalization is innocuous and we adopt it.

The post-search standard deviation \(\tilde\sigma\) is a different
matter. It is identified in principle by functional form ---
\(\tilde\sigma\) appears on both sides of
\(g(\zeta) = \cost/\tilde\sigma\) --- but Yavorsky et al. (2021) shows
by simulation that functional-form identification alone is not enough to
estimate it in practice, and that an exogenous, alternative-specific
cost shifter is what makes estimation work.

The consequence of fixing it deserves to be stated without hedging. With
both variances normalized, search costs cannot be monetized and
search-cost counterfactuals are not interpretable (2 §4.4.4). That would
rule out every friction counterfactual in Section~\ref{sec-managerial}.
Estimating \(\tilde\sigma\) is therefore load-bearing for this paper
rather than a refinement, and it is the single clearest reason SIFT
needs a randomized position rather than merely a rich attribute design.

\subsection{Testing the Zero-Search-Cost
Restriction}\label{sec-est-nesting}

SIFT nests conjoint: as the search cost goes to zero every alternative
is inspected, choice is made over fully realized utilities, and the
model collapses to a standard choice-based conjoint model. Whether
search costs are distinguishable from zero is therefore a testable
question rather than a maintained assumption --- but the test requires
care.

\begin{tcolorbox}[enhanced jigsaw, arc=.35mm, bottomrule=.15mm, bottomtitle=1mm, breakable, colback=white, colbacktitle=quarto-callout-warning-color!10!white, colframe=quarto-callout-warning-color-frame, coltitle=black, left=2mm, leftrule=.75mm, opacityback=0, opacitybacktitle=0.6, rightrule=.15mm, title=\textcolor{quarto-callout-warning-color}{\faExclamationTriangle}\hspace{0.5em}{The obvious test is invalid}, titlerule=0mm, toprule=.15mm, toptitle=1mm]

The natural move is a Wald test on the estimated search cost. It does
not work. Zero is a limit rather than an interior point of the parameter
space: \(\cost \to 0\) corresponds to \(g(\cost) \to +\infty\) and
\(\log \cost \to
-\infty\). The null sits at a boundary under every parameterization, so
the standard asymptotics do not apply.

\end{tcolorbox}

We construct the test from observable implications instead, in three
ways. The first is holdout fit: SIFT against the nested conjoint model
on tasks withheld from estimation. The second is the observable
implication itself --- whether any respondent leaves alternatives
unopened, which under zero search costs no one would. The third is a
likelihood-ratio comparison against a \emph{simulated} rather than
\(\chi^2\) reference distribution.

This also bears on a parameterization choice. 2 §5.2 notes that
estimating \(g\) directly avoids a root-find inside the optimization
loop. Convenient --- but it puts the quantity we most want to test on a
scale where the null is at infinity. We estimate \(\log \cost\) and pay
the root-find, which in our implementation is a spline evaluation rather
than an iterative solve (see
\href{../R/README.md}{\texttt{R/README.md}}).

\subsection{Individual-Level Estimation from a Task
Panel}\label{sec-est-panel}

There are two routes to the individual-level deliverable that makes SIFT
a conjoint substitute rather than a search paper.

\textbf{Full hierarchical Bayes} is what the conjoint audience expects
and what makes individual part-worths and individual search costs
directly reportable.

\textbf{Two-stage empirical Bayes} estimates the hyperparameters by
simulated maximum likelihood and then draws individual posteriors by
Metropolis--Hastings conditional on those estimates. This is what
Morozov et al. (2021) does to compute personalized prices (his §7.2 and
Appendix H), and it is considerably cheaper.

Our runtime work reverses the ordering we originally expected. Because
HB never integrates over the heterogeneity distribution --- each
respondent sits at their own \(\bftheta_i\) --- the mixing-draw
multiplier that makes random-coefficients SMLE expensive simply does not
arise. We therefore treat HB as the primary estimator and empirical
Bayes as the robustness check.

The runtime claim is worth stating precisely, because our first estimate
of it was wrong. Extrapolating from per-evaluation costs we expected
roughly 11 hours for a commercial-scale fit. The implemented sampler
measures at \textbf{1.8 ms per task per iteration} --- linear in tasks,
and only weakly increasing in the number of simulation draws, because
per-task overhead rather than draw count dominates. For 800 respondents
× 10 tasks and 20,000 iterations that is \textbf{about 80 hours
single-threaded}, not 11. The earlier figure assumed one pass over the
data per iteration; the sampler makes two, and carries per-respondent
overhead besides.

What rescues the budget is that the \(\bftheta_i\) step is
embarrassingly parallel: each respondent's Metropolis update touches
only that respondent's tasks. Parallelised across 12 cores we measure a
7.5× speedup --- conservative, since at that problem size communication
is a material share of each iteration --- bringing a commercial fit to
roughly \textbf{11 hours}, inside an overnight budget.

The honest form of the claim is therefore conditional. HB is the cheaper
route to heterogeneity \emph{and} it fits in the budget, \textbf{but
only in parallel}. A single-threaded implementation does not, and any
replication attempt should be told so.

\subsection{Estimation}\label{sec-est-estimation}

Runtime is a binding practical constraint rather than an implementation
detail: a fit has to complete overnight, not over a week. The published
timings make the point. 2 Table 2, for 1,000 consumers with
heterogeneity, report 5h20m for importance sampling, 29h48m for the
kernel-smoothed simulator plus roughly 32 further hours to calibrate its
scaling vector, and 198h24m for GHK. Chung et al. (2024) Table 1, for a
cross-sectional design of comparable size, reports 0.48 minutes for
their probability-mapping simulator against 5--7 minutes kernel-smoothed
and 15--127 minutes for crude frequency --- one to two orders of
magnitude.

We have ported and validated the Chung et al. (2024) estimator and
reproduced their Table 1; the port, its five validation checks, and two
bugs we found in their shipped code are documented in
\href{../R/README.md}{\texttt{R/README.md}}.

For SIFT itself we exploit a structural feature of the design.
Conditional on the observed purchase, every Weitzman condition is a
conjunction of linear inequalities --- the stopping rule is a statement
about a maximum and therefore a union in general, but the choice rule
names which alternative attains that maximum. Continuation then
collapses to the single binding constraint \(\resv_{s_K}\), because
reservation values decrease along the search order. And conditional on
\(\bfeta\) and on the post-search shock to the chosen alternative, the
remaining shocks are independent one-sided constraints that integrate
out in closed form, leaving a one-dimensional integral.

This yields a hybrid estimator: GHK on the pre-search block, quadrature
on the post-search block. Measured against the kernel-smoothed
accept--reject simulator of Yavorsky et al. (2021) at equal draw counts,
its simulation variance is lower by a factor of 21 to 100 depending on
the design. Because runtime is exactly linear in the number of draws,
the relevant summary is time to a given precision, on which the hybrid
is 3 to 21 times cheaper; crude frequency is not competitive on either
margin. The kernel-smoothed simulator also shows the upward bias in the
negative log-likelihood that Yavorsky et al. (2021) notes is intrinsic
to it. The construction and its validation are in
\href{../notebooks/sift_likelihood.qmd}{\texttt{notebooks/sift\_likelihood.qmd}}.

\subsection{Software}\label{sec-est-software}

All estimation code is in R and is released with the paper. The
reservation-value inverse \(g^{-1}\) is evaluated by a cubic spline
built once in the standardized argument \(\cost/\tilde\sigma\), so a
single spline serves every \(\tilde\sigma\); it covers costs down to
\(3.9\times10^{-16}\), well below the range any optimizer will propose.

\section{Simulation Study}\label{sec-simstudy}

\subsection{Design}\label{design}

\subsection{Parameter Recovery}\label{parameter-recovery}

\subsection{Power to Detect a Nonzero Search
Cost}\label{power-to-detect-a-nonzero-search-cost}

\subsection{Misspecification and Behavioral
Robustness}\label{misspecification-and-behavioral-robustness}

\section{Empirical Application}\label{sec-analysis}

\subsection{Data and Design}\label{data-and-design}

\subsection{Descriptive Evidence}\label{descriptive-evidence}

\subsection{Estimates}\label{estimates}

\subsection{Convergent Validity against
Conjoint}\label{convergent-validity-against-conjoint}

\subsection{Holdout Prediction}\label{holdout-prediction}

\subsection{Task-Order Effects}\label{task-order-effects}

\section{What the Model Delivers}\label{sec-managerial}

\subsection{Decomposing Non-Purchase}\label{decomposing-non-purchase}

\subsection{Position and Merchandising
Counterfactuals}\label{position-and-merchandising-counterfactuals}

\subsection{Information Architecture}\label{information-architecture}

\subsection{Search-Based Segmentation}\label{search-based-segmentation}

\section{Discussion}\label{sec-discussion}

\subsection{What the Nesting Result
Changes}\label{what-the-nesting-result-changes}

\subsection{Implications for Practice}\label{implications-for-practice}

\subsection{Limitations}\label{limitations}

\subsection{Where Conjoint Remains the Right
Tool}\label{where-conjoint-remains-the-right-tool}

\section{Conclusion}\label{sec-conclusion}

\section{Appendix}\label{sec-appendix}

\subsection{Notation Crosswalk}\label{sec-appendix-notation}

The two papers this one builds on use different symbols for the same
objects, and in two places the same symbol for different objects. The
table records the mapping, both for readers arriving from either
literature and to keep our own implementation honest against the two
reference codebases.

\begin{longtable}[]{@{}
  >{\raggedright\arraybackslash}p{(\linewidth - 6\tabcolsep) * \real{0.2500}}
  >{\raggedright\arraybackslash}p{(\linewidth - 6\tabcolsep) * \real{0.2500}}
  >{\raggedright\arraybackslash}p{(\linewidth - 6\tabcolsep) * \real{0.2500}}
  >{\raggedright\arraybackslash}p{(\linewidth - 6\tabcolsep) * \real{0.2500}}@{}}
\caption{Notation crosswalk}\label{tbl-notation}\tabularnewline
\toprule\noalign{}
\begin{minipage}[b]{\linewidth}\raggedright
Object
\end{minipage} & \begin{minipage}[b]{\linewidth}\raggedright
This paper
\end{minipage} & \begin{minipage}[b]{\linewidth}\raggedright
2
\end{minipage} & \begin{minipage}[b]{\linewidth}\raggedright
Yavorsky et al. (2021)
\end{minipage} \\
\midrule\noalign{}
\endfirsthead
\toprule\noalign{}
\begin{minipage}[b]{\linewidth}\raggedright
Object
\end{minipage} & \begin{minipage}[b]{\linewidth}\raggedright
This paper
\end{minipage} & \begin{minipage}[b]{\linewidth}\raggedright
2
\end{minipage} & \begin{minipage}[b]{\linewidth}\raggedright
Yavorsky et al. (2021)
\end{minipage} \\
\midrule\noalign{}
\endhead
\bottomrule\noalign{}
\endlastfoot
Utility & \(u_{itj}\) & \(u_{ij}\) & \(u_{ij}\) \\
Pre-search utility & \(\delta_{itj}\) & \(\delta_{ij}\) &
\(\delta_{ij}\) \\
Observed pre-search utility & \(\xi_{itj} = \bfX_{itj}'\bfbeta_i\) &
\(\xi_{ij} = X_j'\bfbeta_i - \alpha_i p_{ij}\) & \(x_j'\bfbeta\) \\
Pre-search taste shock & \(\mu_{itj}\) & \(\mu_{ij}\) & \(\eta_{ij}\) \\
Post-search taste shock & \(\varepsilon_{itj}\) & \(\varepsilon_{ij}\) &
\(\varepsilon_{ij}\) \\
Observed post-search utility & \(\bfL_{itj}'\bfkappa_i\) &
\(L_j'\kappa_i\) (§2.4 only) & --- \\
Post-search s.d. & \(\sigma_\varepsilon\) & \(\sigma_\varepsilon\) &
\(\sigma\) (``MVSD'') \\
Search cost & \(\cost_{itj}\) & \(c_{ij} = Z_j'\gamma_i\) &
\(c_{ij} = \exp\{\gamma_0 + d_{ij}'\gamma\}\) \\
Reservation value & \(\resv_{itj}\) & \(z_{ij}\) & \(z_{ij}\) \\
Reservation offset & \(g(\cost)\) & \(g(c)\),
\(m(c/\sigma_\varepsilon)\) & \(\zeta_{ij}\) \\
Searched set & \(\Sset_{it}\) & \(S_i\) & (implicit) \\
Unsearched set & \(\Sbar_{it}\) & \(\bar S_i\) & (implicit) \\
Number of searches & \(H_{it}\) & \(H_i\) & \(K_i\) \\
Search-order index & \(h\) & \(h\) & (folded into \(j\)) \\
Purchase & \(y_{it}\) & \(y_i\) & \(j^*\) \\
Parameter vector & \(\bftheta_i\) & \(\theta\) & ---
(\(\theta = \log\sigma\)) \\
Population parameters & \(\bfOmega\) & \(\Omega\) & --- \\
Smoothing constants & \(\rho_k\) & \(\rho_k\) & \(\lambda_k\) \\
\end{longtable}

Four collisions are worth naming explicitly, because each is a live
source of error when porting code or reading the two papers side by
side.

\begin{enumerate}
\def\labelenumi{\arabic{enumi}.}
\tightlist
\item
  \textbf{The differenced conditions \(\nu_1\) and \(\nu_2\) are swapped
  between the two papers.} 2 label \(\nu_{1h}\) the selection (order)
  condition and \(\nu_{2h}\) the continuation condition; Yavorsky et al.
  (2021) label \(\nu_{1}\) continuation and \(\nu_{2}\) selection.
  Smoothing constants therefore do not correspond element-wise across
  the two implementations. We follow

  \begin{enumerate}
  \def\labelenumii{\arabic{enumii}.}
  \setcounter{enumii}{1}
  \tightlist
  \item
  \end{enumerate}
\item
  \textbf{\(\theta\) denotes different objects.} 2 use it for the full
  parameter vector; Yavorsky et al. (2021) use it specifically for
  \(\log \sigma\). We follow the former.
\item
  \textbf{\(\zeta\) and \(\xi\) are visually near-identical in print}
  and denote unrelated objects across the two papers --- the
  reservation-value offset in one, observed pre-search utility in the
  other. We retain \(\xi\) and write the offset as \(g(\cdot)\).
\item
  \textbf{\(d\) is a simulation-draw index in one and a distance vector
  in the other.} We use it for neither.
\end{enumerate}

\subsection{Additional Appendices}\label{additional-appendices}

\section*{References}\label{references}
\addcontentsline{toc}{section}{References}

\protect\phantomsection\label{refs}
\begin{CSLReferences}{1}{1}
\bibitem[\citeproctext]{ref-Allenby_2019}
Allenby, Greg M., Nino Hardt, and Peter E. Rossi. 2019. {``Economic
Foundations of Conjoint Analysis.''} In \emph{Handbook of the Economics
of Marketing}. Elsevier.
\url{https://doi.org/10.1016/bs.hem.2019.04.002}.

\bibitem[\citeproctext]{ref-Bronnenberg_2016}
Bronnenberg, Bart J., Jun B. Kim, and Carl F. Mela. 2016. {``Zooming in
on Choice: How Do Consumers Search for Cameras Online?''}
\emph{Marketing Science} 35 (5): 693--712.
\url{https://doi.org/10.1287/mksc.2016.0977}.

\bibitem[\citeproctext]{ref-Brown_2011}
Brown, Meta, Christopher J Flinn, and Andrew Schotter. 2011.
{``Real-Time Search in the Laboratory and the Market.''} \emph{American
Economic Review} 101 (2): 948--74.
\url{https://doi.org/10.1257/aer.101.2.948}.

\bibitem[\citeproctext]{ref-Caplin_2011}
Caplin, Andrew, Mark Dean, and Daniel Martin. 2011. {``Search and
Satisficing.''} \emph{American Economic Review} 101 (7): 2899--922.
\url{https://doi.org/10.1257/aer.101.7.2899}.

\bibitem[\citeproctext]{ref-Chen_Yao_2017}
Chen, Yuxin, and Song Yao. 2017. {``Sequential Search with Refinement:
Model and Application with Click-Stream Data.''} \emph{Management
Science} 63 (12): 4345--65.
\url{https://doi.org/10.1287/mnsc.2016.2557}.

\bibitem[\citeproctext]{ref-Chung_2025}
Chung, Jae Hyen, Pradeep Chintagunta, and Sanjog Misra. 2024.
{``Simulated Maximum Likelihood Estimation of the Sequential Search
Model.''} \emph{Quantitative Marketing and Economics} 23 (1): 105--64.
\url{https://doi.org/10.1007/s11129-024-09281-4}.

\bibitem[\citeproctext]{ref-Cox_Oaxaca_1989}
Cox, JamesC., and RonaldL. Oaxaca. 1989. {``Laboratory Experiments with
a Finite-Horizon Job-Search Model.''} \emph{Journal of Risk and
Uncertainty} 2 (3). \url{https://doi.org/10.1007/bf00209391}.

\bibitem[\citeproctext]{ref-DeLosSantos_2012}
De los Santos, Babur, Ali Hortaçsu, and Matthijs R Wildenbeest. 2012.
{``Testing Models of Consumer Search Using Data on Web Browsing and
Purchasing Behavior.''} \emph{American Economic Review} 102 (6):
2955--80. \url{https://doi.org/10.1257/aer.102.6.2955}.

\bibitem[\citeproctext]{ref-Ding_2007}
Ding, Min. 2007. {``An Incentive-Aligned Mechanism for Conjoint
Analysis.''} \emph{Journal of Marketing Research} 44 (2): 214--23.
\url{https://doi.org/10.1509/jmkr.44.2.214}.

\bibitem[\citeproctext]{ref-Ding_2005}
Ding, Min, Rajdeep Grewal, and John Liechty. 2005. {``Incentive-Aligned
Conjoint Analysis.''} \emph{Journal of Marketing Research} 42 (1):
67--82. \url{https://doi.org/10.1509/jmkr.42.1.67.56890}.

\bibitem[\citeproctext]{ref-Dzyabura_2011}
Dzyabura, Daria, and John R. Hauser. 2011. {``Active Machine Learning
for Consideration Heuristics.''} \emph{Marketing Science} 30 (5):
801--19. \url{https://doi.org/10.1287/mksc.1110.0660}.

\bibitem[\citeproctext]{ref-Gabaix_2006}
Gabaix, Xavier, David Laibson, Guillermo Moloche, and Stephen Weinberg.
2006. {``Costly Information Acquisition: Experimental Analysis of a
Boundedly Rational Model.''} \emph{American Economic Review} 96 (4):
1043--68. \url{https://doi.org/10.1257/aer.96.4.1043}.

\bibitem[\citeproctext]{ref-Gilbride_Allenby_2004}
Gilbride, Timothy J., and Greg M. Allenby. 2004. {``A Choice Model with
Conjunctive, Disjunctive, and Compensatory Screening Rules.''}
\emph{Marketing Science} 23 (3): 391--406.
\url{https://doi.org/10.1287/mksc.1030.0032}.

\bibitem[\citeproctext]{ref-Greminger_2022}
Greminger, Rafael P. 2022. {``Optimal Search and Discovery.''}
\emph{Management Science} 68 (5): 3904--24.
\url{https://doi.org/10.1287/mnsc.2021.4085}.

\bibitem[\citeproctext]{ref-Hauser_2014}
Hauser, John R. 2014. {``Consideration-Set Heuristics.''} \emph{Journal
of Business Research} 67 (8): 1688--99.
\url{https://doi.org/10.1016/j.jbusres.2014.02.015}.

\bibitem[\citeproctext]{ref-Hauser_2010}
Hauser, John R., Olivier Toubia, Theodoros Evgeniou, Rene Befurt, and
Daria Dzyabura. 2010. {``Disjunctions of Conjunctions, Cognitive
Simplicity, and Consideration Sets.''} \emph{Journal of Marketing
Research} 47 (3): 485--96. \url{https://doi.org/10.1509/jmkr.47.3.485}.

\bibitem[\citeproctext]{ref-Hey_1987}
Hey, John D. 1987. {``Still Searching.''} \emph{Journal of Economic
Behavior \&Amp; Organization} 8 (1): 137--44.
\url{https://doi.org/10.1016/0167-2681(87)90026-6}.

\bibitem[\citeproctext]{ref-Hong_Shum_2006}
Hong, Han, and Matthew Shum. 2006. {``Using Price Distributions to
Estimate Search Costs.''} \emph{The RAND Journal of Economics} 37 (2):
257--75. \url{https://doi.org/10.1111/j.1756-2171.2006.tb00015.x}.

\bibitem[\citeproctext]{ref-Honka_2014}
Honka, Elisabeth. 2014. {``Quantifying Search and Switching Costs in the
{US} Auto Insurance Industry.''} \emph{The RAND Journal of Economics} 45
(4): 847--84. \url{https://doi.org/10.1111/1756-2171.12073}.

\bibitem[\citeproctext]{ref-Honka_Chintagunta_2017}
Honka, Elisabeth, and Pradeep Chintagunta. 2017. {``Simultaneous or
Sequential? Search Strategies in the u.s. Auto Insurance Industry.''}
\emph{Marketing Science} 36 (1): 21--42.
\url{https://doi.org/10.1287/mksc.2016.0995}.

\bibitem[\citeproctext]{ref-Hortacsu_Syverson_2004}
Hortacsu, A., and C. Syverson. 2004. {``Product Differentiation, Search
Costs, and Competition in the Mutual Fund Industry: A Case Study of
s\&amp;p 500 Index Funds.''} \emph{The Quarterly Journal of Economics}
119 (2): 403--56. \url{https://doi.org/10.1162/0033553041382184}.

\bibitem[\citeproctext]{ref-Kim_2010}
Kim, Jun B., Paulo Albuquerque, and Bart J. Bronnenberg. 2010. {``Online
Demand Under Limited Consumer Search.''} \emph{Marketing Science} 29
(6): 1001--23. \url{https://doi.org/10.1287/mksc.1100.0574}.

\bibitem[\citeproctext]{ref-Koulayev_2014}
Koulayev, Sergei. 2014. {``Search for Differentiated Products:
Identification and Estimation.''} \emph{The RAND Journal of Economics}
45 (3): 553--75. \url{https://doi.org/10.1111/1756-2171.12062}.

\bibitem[\citeproctext]{ref-Marshall_Bradlow_2002}
Marshall, Pablo, and Eric T Bradlow. 2002. {``A Unified Approach to
Conjoint Analysis Models.''} \emph{Journal of the American Statistical
Association} 97 (459): 674--82.
\url{https://doi.org/10.1198/016214502388618410}.

\bibitem[\citeproctext]{ref-McCall_1970}
McCall, J. J. 1970. {``Economics of Information and Job Search.''}
\emph{The Quarterly Journal of Economics} 84 (1): 113.
\url{https://doi.org/10.2307/1879403}.

\bibitem[\citeproctext]{ref-Meissner_2016}
Meißner, Martin, Andres Musalem, and Joel Huber. 2016. {``Eye Tracking
Reveals Processes That Enable Conjoint Choices to Become Increasingly
Efficient with Practice.''} \emph{Journal of Marketing Research} 53 (1):
1--17. \url{https://doi.org/10.1509/jmr.13.0467}.

\bibitem[\citeproctext]{ref-Morozov_2021}
Morozov, Ilya, Stephan Seiler, Xiaojing Dong, and Liwen Hou. 2021.
{``Estimation of Preference Heterogeneity in Markets with Costly
Search.''} \emph{Marketing Science} 40 (5): 871--99.
\url{https://doi.org/10.1287/mksc.2021.1287}.

\bibitem[\citeproctext]{ref-Payne_1988}
Payne, John W., James R. Bettman, and Eric J. Johnson. 1988. {``Adaptive
Strategy Selection in Decision Making.''} \emph{Journal of Experimental
Psychology: Learning, Memory, and Cognition} 14 (3): 534--52.
\url{https://doi.org/10.1037/0278-7393.14.3.534}.

\bibitem[\citeproctext]{ref-Payne_1993}
Payne, John W., James R. Bettman, and Eric J. Johnson. 1993. \emph{The
Adaptive Decision Maker}. Cambridge University Press.
\url{https://doi.org/10.1017/cbo9781139173933}.

\bibitem[\citeproctext]{ref-Reutskaja_2011}
Reutskaja, Elena, Rosemarie Nagel, Colin F Camerer, and Antonio Rangel.
2011. {``Search Dynamics in Consumer Choice Under Time Pressure: An
Eye-Tracking Study.''} \emph{American Economic Review} 101 (2):
900--926. \url{https://doi.org/10.1257/aer.101.2.900}.

\bibitem[\citeproctext]{ref-Schotter_Braunstein_1981}
SCHOTTER, ANDREW, and YALE M. BRAUNSTEIN. 1981. {``Economic Search: An
Experimental Study.''} \emph{Economic Inquiry} 19 (1): 1--25.
\url{https://doi.org/10.1111/j.1465-7295.1981.tb00600.x}.

\bibitem[\citeproctext]{ref-Sonnemans_1998}
Sonnemans, Joep. 1998. {``Strategies of Search.''} \emph{Journal of
Economic Behavior \&Amp; Organization} 35 (3): 309--32.
\url{https://doi.org/10.1016/s0167-2681(98)00051-1}.

\bibitem[\citeproctext]{ref-Stigler_1961}
Stigler, George J. 1961. {``The Economics of Information.''}
\emph{Journal of Political Economy} 69 (3): 213--25.
\url{https://doi.org/10.1086/258464}.

\bibitem[\citeproctext]{ref-Ursu_2018}
Ursu, Raluca M. 2018. {``The Power of Rankings: Quantifying the Effect
of Rankings on Online Consumer Search and Purchase Decisions.''}
\emph{Marketing Science} 37 (4): 530--52.
\url{https://doi.org/10.1287/mksc.2017.1072}.

\bibitem[\citeproctext]{ref-Ursu_2020}
Ursu, Raluca M., Qingliang Wang, and Pradeep K. Chintagunta. 2020.
{``Search Duration.''} \emph{Marketing Science} 39 (5): 849--71.
\url{https://doi.org/10.1287/mksc.2020.1225}.

\bibitem[\citeproctext]{ref-Ursu_Seiler_Honka_2024}
Ursu, Raluca, Stephan Seiler, and Elisabeth Honka. 2024. {``The
Sequential Search Model: A Framework for Empirical Research.''}
\emph{Quantitative Marketing and Economics} 23 (1): 165--213.
\url{https://doi.org/10.1007/s11129-024-09291-2}.

\bibitem[\citeproctext]{ref-Weitzman_1979}
Weitzman, Martin L. 1979. {``Optimal Search for the Best Alternative.''}
\emph{Econometrica} 47 (3): 641--54.
\url{https://doi.org/10.2307/1910412}.

\bibitem[\citeproctext]{ref-Yang_2015}
Yang, Liu (Cathy), Olivier Toubia, and Martijn G. De Jong. 2015. {``A
Bounded Rationality Model of Information Search and Choice in Preference
Measurement.''} \emph{Journal of Marketing Research} 52 (2): 166--83.
\url{https://doi.org/10.1509/jmr.13.0288}.

\bibitem[\citeproctext]{ref-Yang_2018}
Yang, Liu (Cathy), Olivier Toubia, and Martijn G. de Jong. 2018.
{``Attention, Information Processing, and Choice in Incentive-Aligned
Choice Experiments.''} \emph{Journal of Marketing Research}, ahead of
print. \url{https://doi.org/10.1509/jmr.15.0474}.

\bibitem[\citeproctext]{ref-Yavorsky_2021}
Yavorsky, Dan, Elisabeth Honka, and Keith Chen. 2021. {``Consumer Search
in the {U.S.} Auto Industry: The Role of Dealership Visits.''}
\emph{Quantitative Marketing and Economics} 19 (1): 1--52.
\url{https://doi.org/10.1007/s11129-020-09229-4}.

\bibitem[\citeproctext]{ref-Yee_2007}
Yee, Michael, Ely Dahan, John R. Hauser, and James Orlin. 2007.
{``Greedoid-Based Noncompensatory Inference.''} \emph{Marketing Science}
26 (4): 532--49. \url{https://doi.org/10.1287/mksc.1060.0213}.

\end{CSLReferences}




\end{document}
